% 03_macroscopic_relativity.tex
\chapter{Macroscopic Relativity: The Optical Metric}
\label{ch:relativity}


\begin{objectivebox}
\begin{itemize}
    \item Formulate General Relativity as Applied Vacuum Engineering: the optical metric of a refractive dielectric medium.
    \item Derive the Schwarzschild metric as the impedance profile of a spherically symmetric vacuum compression.
    % Two longitudinal speeds: solid P-wave $c_L=\sqrt{10/3}\,c$ (carries the $4G/3$ shear term); bulk-dilatational $\sqrt{2}\,c$ drops the shear term.
    \item Understand why $c$ is the maximum \emph{transverse} signal speed but the longitudinal (P) wave propagates at the isotropic-solid speed $c_L = \sqrt{(K_{vac} + \tfrac{4}{3}G_{vac})/\rho_{bulk}} = \sqrt{\tfrac{10}{3}}\,c \approx 1.83c$ at the $K_{vac} = 2G_{vac}$ operating point (equivalently Poisson ratio $\nu = 2/7$; canonical derivation in Vol~2 Ch~7). The pure bulk-modulus speed $\sqrt{K_{vac}/\rho_{bulk}} = \sqrt{2}\,c$ is the dilatational compression speed that drops the $4G/3$ shear term, not the P-wave. \textbf{[DEMOTED 2026-08-11 --- R40-B1; dated demotion note at the end of this chapter]}
\end{itemize}
\end{objectivebox}

Standard pedagogical models of General Relativity often rely on the heuristic of a 2D elastic membrane warping into an additional spatial dimension. The AVE framework offers an alternative formulation grounded in the Electrodynamics of a \textbf{3D Optical Impedance Metric}.

\section{Gravity as Applied Vacuum Engineering}
In the AVE framework, the macroscopic effective vacuum is modelled as a 3D Electromagnetic LC Network. When a massive topological defect (a confined light knot or star) forms, its localised inductive rest-energy polarises the surrounding spatial discrete edges. This polarisation \textbf{compresses the effective impedance} ($\epsilon\mu$) inward toward the centre of mass.

Geometrically polarizing these edges into a smaller volume locally increases the absolute optical density of the spatial substrate, yielding a proportional increase in the localized \textbf{Refractive Index ($n$)}. Gravitational attraction is thus modeled entirely via the \textbf{Ponderomotive Force}. A wave packet refracts toward the region of highest dielectric density: its phase velocity $c/n(r)$ is slower where $n$ is larger, so the graded-index medium bends the ray down the impedance gradient (Fermat / Snell refraction), exactly as a wave steers in any graded-index dielectric --- no energy-basin descent is invoked. Gravity represents the refraction of physical confined light down a 3D dielectric impedance gradient.

\subsection{Deriving the Refractive Gradient from LC Polarization}

The macroscopic vacuum moduli are elevated from scalars to rank-2 symmetric tensors. As
established historically by the Gordon Optical Metric, signal propagation through an
anisotropic continuous dielectric reproduces geodesic paths in curved spacetime:

\begin{resultbox}{Gordon Optical Metric}
\begin{equation}
g_{\mu\nu}^{AVE}=\eta_{\mu\nu}+\left(1-\frac{1}{n^{2}(r)}\right)u_{\mu}u_{\nu}
\end{equation}
\end{resultbox}

By applying standard 3D electrostatics using the Laplace equation against a steady-state inductive energy density ($M$), balanced against the continuous macroscopic electrodynamic impedance limit ($T_{max,g} = \xi T_{EM} = c^{4}/7G$), the localised \textbf{1D principal radial polarisation strain} ($\epsilon_{11}$) field generates the $1/r$ Newtonian potential:

\begin{resultbox}{Radial Polarisation Strain Poisson Equation}
\begin{equation}
-\left(\frac{c^{4}}{7G}\right)\nabla^{2}\epsilon_{11}(r)=4\pi Mc^{2}\delta^{3}(r)
\end{equation}
\end{resultbox}

Convolving this source with the 3D Laplacian Green's function ($-1/4\pi r$) yields the steady-state 1D principal radial strain field:

\begin{resultbox}{Radial Strain to Newtonian Potential}
\begin{equation}
\epsilon_{11}(r)=\frac{7GM}{c^{2}r}
\end{equation}
\end{resultbox}

\section{The Ponderomotive Equivalence Principle}

Standard physics invokes the Weak Equivalence Principle ($m_i=m_g$) as an axiomatic postulate. AVE derives it from macroscopic wave mechanics.

Because a massive topological wave-packet acts as a 3D isotropic defect, it couples to the spatial volume via the $1/7$ Lagrangian isotropic projection (derived in Chapter 1). The effective scalar refractive index perceived by mass is $n_{scalar}(r)=1+\epsilon_{11}(r)/7=1+GM/c^{2}r$. The localised stored energy of the knot is its internal inductive rest mass ($m_i c^{2}$) scaled inversely by the refractive density:

\begin{resultbox}{Ponderomotive Wave Energy}
\begin{equation}
U_{wave}(r)=\frac{m_i c^{2}}{n_{scalar}(r)}\approx m_i c^{2}\left(1-\frac{GM}{rc^{2}}\right)=m_i c^{2}-\frac{GM m_i}{r}
\end{equation}
\end{resultbox}

Taking the spatial gradient directly yields the gravitational acceleration, expressed as $F_{grav}=-\nabla U_{wave}$:

\begin{resultbox}{Newtonian Gravity from Optical Gradient}
\begin{equation}
F_{grav}=-\frac{GM m_i}{r^{2}}\hat{r}
\end{equation}
\end{resultbox}

Because the localised wave energy is defined by the particle's inductive inertia $m_i$, it cancels out of the acceleration equation ($F=ma$), guaranteeing that inertial mass and gravitational mass are identical ($m_i\equiv m_g$).

\section{The Optical Metric: Gravity as Refractive Density}
\label{sec:optical_metric}

Standard General Relativity models gravity as coordinate curvature. In the AVE framework, gravity is defined as the \textbf{Electromagnetic Densification} of the vacuum LC network. A massive object acts as a refractive index sink, polarising the surrounding network density.

\subsection{Deriving the Refractive Index}

The macroscopic vacuum moduli are elevated to rank-2 symmetric tensors. Because the vacuum acts macroscopically as a Trace-Reversed Continuum to support transverse EM waves, it possesses a fixed effective Poisson ratio of $\nu_{vac}=2/7$. 

A localised massive defect does not exert a uniform 3D hydrostatic compression; it exerts a radial pull, acting as the continuous source of the principal radial tensile strain ($\epsilon_{11} > 0$). Conversely, light propagates as a transverse shear wave and couples to the orthogonal transverse spatial metric ($h_{\perp}$). 

In continuum mechanics, radial tension causes orthogonal transverse space to contract ($h_{\perp} = -\nu_{vac} \cdot \epsilon_{11}$). However, the effective refractive index ($n$) scales proportionally with the physical geometric optical density ($\rho_{opt}$) of the medium. Because optical density scales inversely with physical transverse spatial displacement ($n \propto \frac{1}{1 + h_{\perp}}$), applying the first-order Taylor expansion for small macroscopic strains ($\frac{1}{1+x} \approx 1-x$), a negative (compressive) transverse strain yields a positive increase in the effective refractive index:

\begin{resultbox}{Transverse Refractive Index}
\begin{equation}
n(r) = 1 - h_{\perp} = 1 - (-\nu_{vac}\epsilon_{11}) = 1 + \nu_{vac}\epsilon_{11}
\end{equation}
\end{resultbox}

Substituting the trace-reversed tensor boundary ($\nu_{vac} = 2/7$) and the radial strain field yields:

\begin{resultbox}{The Refractive Index of Gravity}
\begin{equation}
n(r) = 1 + \left(\frac{2}{7}\right)\left(\frac{7GM}{c^2r}\right) = 1 + \frac{2GM}{c^2r}
\end{equation}
\end{resultbox}

The effective Refractive Index ($n$) perceived by a photon is therefore mathematically identical to the spatial transverse trace of the Gordon optical metric.

\begin{figure}[h!]
    \centering
    \includegraphics[width=0.85\textwidth]{optical_refractive_gradient.png}
    \caption{\textbf{Gravity as Optical Dielectric Refraction.} The presence of a localized topological defect (Mass) polarizes the surrounding continuous dielectric LC grid. This reactive detuning raises the local refractive index ($n > 1$) while holding the characteristic impedance invariant ($Z = Z_0$, reflectionless --- Symmetric Gravity). In the AVE framework, "Gravity" is structurally a frequency-detuned (slower-ringing) refractive gradient at invariant $Z_0$ --- not a denser/higher-impedance medium; the diffraction of light into it is reflectionless precisely because $Z$ does not change.}
    \label{fig:optical_refractive_gradient}
\end{figure}

\subsection{The Absolute Intergalactic Speed of Light ($c_{max}$)}
Because the physical speed of light ($c_{local}$) is governed inversely by the local LC refractive index ($c_{local} = c_0/n$), General Relativity's assertion that $c$ is a rigid, universal constant evaluated identically everywhere does not hold within the AVE framework.

Earth resides within a compound gravitational well generated by the localised mass of the Earth, the Sun, and the Milky Way galaxy. The lowest theoretical density of the LC network occurs in the deepest voids of intergalactic space, where the ambient gravitational potential ($\Phi \rightarrow 0$) approaches zero. 

By calculating the total fractional change in the refractive index induced upon the local solar system by the Milky Way's galactic tensor strain ($\Phi_{galaxy} = GM/R \approx 5.4 \times 10^{11} \text{ J/kg}$):
\begin{resultbox}{Galactic Refractive Deficit}
\begin{equation}
    \frac{\Delta c}{c_{local}} = \frac{2\Phi_{galaxy}}{c^2} = \frac{2(5.4 \times 10^{11})}{(3 \times 10^8)^2} \approx 1.2 \times 10^{-5}
\end{equation}
\end{resultbox}

This ratio ($1{:}1.000012$) dictates that the local speed of light measured on Earth ($299{,}792{,}458$ m/s) is artificially constrained by ambient galactic dielectric density. In the undisturbed, fully relaxed state of intergalactic space, the absolute unconstrained maximum speed of light accelerates by approximately ${\sim}\,3{,}600$ m/s to $c_{max} \approx 299{,}796{,}055$ m/s.

This establishes, as an illustrative extrapolation of the same local-refraction relation to $\Phi \to 0$ (not an experimentally validated prediction), the baseline that macroscopic metric engineering (``Warp Transit'') would target: a vessel lowering its local refractive index below the Earth baseline toward superluminal effective propagation ($n \ll 1.0$).

\subsection{Achromatic Impedance Matching}
\label{sec:achromatic_matching}
A critical property of astrophysical gravity is that it behaves as a transparent lens. If gravity is an optical dense metric where the speed of light decreases locally ($c' = 1/\sqrt{\epsilon' \mu'} = c/n(r)$), classical optics predicts that light should suffer partial reflection when crossing an impedance gradient ($Z_1 \neq Z_2$). 

In the AVE framework, this is resolved because the geometric polarization of the LC network scales its dual reactive components symmetrically. The absolute values of local magnetic permeability ($\mu$) and dielectric permittivity ($\epsilon$) both scale directly and proportionately with the local scalar strain: 
\begin{resultbox}{Symmetric Impedance Scaling}
\begin{equation}
    \mu' = n(r)\mu_0 \quad \text{and} \quad \epsilon' = n(r)\epsilon_0
\end{equation}
\end{resultbox}

Consequently, while the local phase velocity is reduced ($c' = 1/\sqrt{n^2\mu_0 \epsilon_0} = c/n(r)$), the local characteristic transverse impedance of the vacuum \textbf{remains strictly invariant}:
\begin{resultbox}{Achromatic Impedance Matching}
\begin{equation}
    Z_0' = \sqrt{\frac{\mu'}{\epsilon'}} = \sqrt{\frac{n(r)\mu_0}{n(r)\epsilon_0}} = \sqrt{\frac{\mu_0}{\epsilon_0}} \equiv Z_0 \approx 376.73\ \Omega
\end{equation}
\end{resultbox}

Because the transverse impedance ratio is preserved across all gravitational gradients, the spatial vacuum operates as an \textbf{Achromatic Impedance-Matched Lens}. This guarantees that propagating transverse light diffracts and bends through deep gravity wells without suffering chromatic dispersion, internal scattering, or boundary back-reflection.


\subsection{Verification: The Einstein Lensing Deflection}
To test this Optical Metric, a numerical ray-tracing simulation integrates a photon trajectory past the Sun. Integrating Snell's Law through this specific refractive gradient yields a total deflection angle of:
\begin{resultbox}{Einstein Lensing Deflection}
\begin{equation}
    \delta = \frac{4GM}{b c^2}
\end{equation}
\end{resultbox}
This result reproduces the Einstein prediction, distinguishing the AVE framework from Newtonian corpuscular models ($\delta = 2GM/bc^2$) without invoking higher-dimensional curvature. This reproduction of the GR weak-field deflection is a consistency check via Gordon optical-metric isomorphism; the AVE-distinct, forward-testable departure appears only in the trans-Planckian discrete-lattice $(q\ell_{node})^4$ dispersion regime.

\section{Resolving the Cauchy Implosion Paradox}

Standard 19th-century aether models were challenged by the Cauchy Implosion Paradox:
enforcing purely transverse wave limits required a negative bulk modulus ($K_{cauchy}=
-4/3G_{vac}$), implying the universe would thermodynamically implode.

The substrate resolves this via its fundamental \textbf{LC Network Topology}. As structurally established
in Chapter 1, the trace-reversed equilibrium of the non-affine substrate rigidly
locks the macroscopic bulk impedance to double the transverse shear impedance ($K_{vac}\equiv2G_{vac}$).
This positive bulk resistance guarantees that the spatial substrate is \textbf{[DEMOTED 2026-08-11 --- R40-B2a: NEEDS RE-DERIVATION, not dead; dated note at the end of this file]}
incompressible and thermodynamically stable against gravitational collapse.

\section{The Event Horizon as Dielectric Rupture}
\label{sec:dielectric_rupture}

The Event Horizon is classically defined as a coordinate singularity. In the AVE framework it resolves into two distinct channel boundaries. As matter aggregates, the local \textbf{EM/transverse} refractive strain ($n(r)-1 = 2GM/(c^2 r)$) increases; it reaches unity at the Schwarzschild radius, the radius at which \emph{light} sees the General Relativistic event horizon ($n \to \infty$). Under Symmetric Gravity the EM impedance is invariant ($Z_{EM}=Z_0$, $\Gamma_{EM}=0$): the EM channel stays matched and does not reflect or rupture there. The $\Gamma=-1$\gammaundeclared{} total-reflection wall is a separate, deeper boundary in the \textbf{shear and bulk} channels, where $\varepsilon_{11}=1$, $G_{shear}\to0$ and the bulk and shear impedances collapse ($Z_{shear},Z_{bulk}\to0$).
\begin{resultbox}{Two channel boundaries (EM horizon vs shear/bulk rupture)}
\begin{equation}
\begin{aligned}
    \text{EM horizon:} &\quad \frac{2GM}{c^2 R_{rupture}} \equiv 1.0 \implies R_{rupture} = \frac{2GM}{c^2}, \\
    \text{shear/bulk rupture:} &\quad r_{sat} = \frac{7GM}{c^2} = 3.5\,r_s
\end{aligned}
\end{equation}
\end{resultbox}
This identifies the Schwarzschild radius not as a point of infinite curvature, but as the physical radius where the EM/transverse refractive strain reaches unity (the optical horizon, impedance-matched $\Gamma_{EM}=0$); the actual $\Gamma=-1$\gammaundeclared{} impedance rupture is the shear and bulk boundary at $r_{sat}=7GM/c^2=3.5\,r_s$.

\section{Gravitomagnetism: Frame Dragging as Mutual Inductance}
\label{sec:gravitomagnetism_viscosity}
In standard General Relativity, a rotating massive body drags the geometric fabric of spacetime along with it—a phenomenon known as the Lense-Thirring effect (Gravitomagnetism). In the AVE framework, because the physical vacuum operates dynamically as an LC network (the substrate), this effect is analytically identical to classical \textbf{Macroscopic Mutual Inductance}.

As a massive macroscopic boundary (a spinning planet) rotates, its circulating current ($I_{mass}$) couples to the adjacent spatial metric layer via \textbf{macroscopic mutual inductance} ($M_{12}$). Macroscopic inductive momentum transport continues across the lattice, inducing a localised, steady-state magnetic bias field in the surrounding vacuum. For a 2D equatorial slice, the steady-state induced rotational vector potential of the LC lattice decays as an inverse square ($\Omega_{\text{induced}} \propto 1/r^2$). This matches the weak-field General Relativistic prediction ($\Omega_{LT}$), deriving Gravitomagnetism as standard macroscopic mutual inductance without invoking additional geometrical abstractions.

Furthermore, this mutual inductance biases the surrounding vacuum into an angle-dependent, non-reciprocal impedance gradient ($n(r,\theta)$). The physical mechanism is this \textbf{impedance-gradient convolution / mutual inductance}: the rotating boundary convolutes the static refractive index into a direction-dependent bias, so a transverse probe's phase velocity shifts as $c' = c/n(r,\theta)$. This direction-dependent phase-velocity shift is the physical cause of the asymmetric D-shaped photon-capture shadows observed around rapidly rotating sources (such as Kerr black holes). \emph{Bounded analogy (not the mechanism):} the differential phase velocity \emph{looks like} a differential-speed flow, inviting a Bernoulli-style intuition ($c' = c/n \pm \vec{v}$), but \textbf{nothing literally advects} --- the bulk substrate is at rest in the CMB frame; the ``velocity field'' is an effective impedance-gradient drift, not an advecting fluid.

\begin{figure}[h]
    \centering
    \includegraphics[width=1.0\textwidth]{lense_thirring_inductive_drag.png}
    \caption{\textbf{Gravitomagnetism as Mutual Inductance.} The Lense-Thirring frame-dragging effect is modelled as the mutual inductance of the substrate LC network. A rotating massive body inductively biases the adjacent vacuum lattice; due to the structural macroscopic inductance ($\mu_{vac}$), angular momentum diffuses radially outward, forming a steady-state magnetic bias field that replicates the General Relativistic $1/r^2$ decay field.}
    \label{fig:lense_thirring_inductive_drag}
\end{figure}

\begin{summarybox}
\begin{itemize}
    \item The Schwarzschild radius is the EM/transverse optical horizon where the refractive index diverges ($n \to \infty$) yet the EM impedance stays matched ($Z_{EM}=Z_0$, $\Gamma_{EM}=0$, no EM reflection); it is the radius at which light sees the General Relativistic horizon.
    \item The $\Gamma=-1$\gammaundeclared{} shear and bulk impedance rupture ($G_{shear}\to0$, $Z_{shear},Z_{bulk}\to0$) is the separate, deeper boundary at $r_{sat}=7GM/c^2=3.5\,r_s$, where the continuous tensor strain reaches the Axiom 4 saturation limit and the spatial mesh undergoes a pre-geometric thermodynamic phase transition.
    \item Gravitomagnetism and the Lense-Thirring effect are rigorously modeled as macroscopic mutual inductance ($M_{12}$) across the substrate.
    \item A rotating massive body couples inductively to the surrounding spatial metric layer, forcing an inverse-square decaying rotational vector potential structurally identical to relativistic frame dragging.
\end{itemize}
\end{summarybox}

\begin{exercisebox}
\begin{enumerate}
    \item Derive the Event Horizon radius specifically identifying the point where unitary dielectric strain breaks the Axiom 4 structural continuity logic of the lattice.
    \item Formulate Gravitomagnetism explicitly as the macroscopic mutual inductance of the LC substrate. Prove that the steady-state induced rotational vector potential mimics the classic Lense-Thirring decay ($\propto 1/r^2$).
\end{enumerate}
\end{exercisebox}

\medskip
\noindent\textbf{[Dated demotion note --- 2026-08-11 (R40 demotion sweep, batch~1;
class DIES-WITH-THE-PHANTOM).]} Status change only: the claim text is \emph{preserved
verbatim} (honesty-lag pattern, Rule~12) and stamped in place; it is no longer live canon.
Nothing is deleted.

\begin{itemize}
  \item \textbf{Line~11} (the objectives box) --- asserts the propagating P-wave outright at $\sqrt{10/3}\,c \approx 1.83c$, plus $\sqrt2\,c$ as ``the dilatational compression speed''. Stamped in place at line~11. It also still carries the stale ``canonical derivation in Vol~2 Ch~7'' provenance for $\nu = 2/7$ that Vol~1 Ch~4 has already corrected to GR-imported (comment twin at line~10) --- that provenance repair is \emph{not} made here; it is a batch-2 re-derivation item.
\end{itemize}

\noindent\textbf{The arc, complete --- the framing R40 rules every demotion note
carries.} (1)~\textbf{The kill fired} (\#930), the walk-back that closed the bulk
radiative-port reading. (2)~\textbf{The premise localized to the \#261 $K = 2G$ import}
(G-RECON, unchallenged): the compressible far-field branch was minted by a GR-imported
elastic modulus, not forced by the axioms. (3)~\textbf{The axioms underdetermine the bulk
sector} --- the \#935 flat-direction finding: the written action conserves the Gauss
function pointwise and never fixes its value. (4)~\textbf{The replacement is the RATIFIED
bound-sector law --- Axiom~5, Substrate DC Bias} (BC-SRC clauses S / G / Q), ratified per
\texttt{\_orchestration/docket-entries/2026-08-10-ruling-r43-ratification.md}, as reconciled by \texttt{\_orchestration/docket-entries/2026-08-10-ruling-r44-r43-reconciliation.md} (R44: the
full-scope R43 record is FINAL and authoritative; the partial
\texttt{2026-08-10-ruling-r43-sg-ratified.md} is SUPERSEDED and is \emph{not} the
resolution). Under the ratified law the A1 / bulk slot is a \textbf{bound response} ---
mechanism gloss \textbf{back-reaction} --- with no independent propagating branch, no port,
and zero longitudinal characteristic speed. A bulk \emph{wave speed}, a bulk \emph{radiative
port}, a bulk \emph{band-branch} and a bulk \emph{transit clock} therefore have no referent.

\noindent\textbf{Standing named-open debt (the honest rider).} The ratified axiom does
\emph{not} discharge everything: \textbf{THE BIAS PROPAGATION THEOREM} is Axiom~5's standing
named-open entry --- clause~G's elliptic law is the \emph{static abstraction} of underived
finite-speed bias dynamics (\texttt{\_orchestration/2026-08-10\_bias-propagation-brief.md}). Where a demoted claim's
replacement depends on finite-speed bias dynamics, the resolution is the ratified axiom
\emph{with that debt open}, not a closed replacement.

\noindent\textbf{Records.} R40 ruling \texttt{\_orchestration/docket-entries/2026-08-10-rulings-r40-r42.md}; verified worklist
\texttt{research/drivers/r40\_sweep\_worklist\_verified.json}; scope verification \texttt{\_orchestration/2026-08-10\_r40-sweep-scope-verification.md};
batch-1 record \texttt{\_orchestration/2026-08-11\_r40-sweep-batch1.md}; vocabulary ruling R50
\texttt{\_orchestration/docket-entries/2026-08-10-ruling-r50-vocab.md} (canonical: the displacement pattern $u_0$ around a deposit is
\textbf{the bound response}, mechanism gloss \textbf{back-reaction}; $\varepsilon_{11}$ is
\textbf{the bias}; ``dress'', ``grade''-as-canonical-noun and ``halo''-for-the-physics are
retired, and the owed theorem is renamed \textbf{THE BIAS PROPAGATION THEOREM}); vocabulary ruling
\textbf{R49(b)} \texttt{\_orchestration/docket-entries/2026-08-10-rulings-r48-r49.md}
(``retardation'' is retired; the canonical term is \textbf{propagation delay / finite propagation
speed} --- the retardation retirement is R49(b)'s, \emph{not} R50's; corrected 2026-08-11 at review).

% --------------------------------------------------------------------------
% R40 batch-1 demotion note --- per-line detail (verbatim audited rows).
% --------------------------------------------------------------------------
%   :11  QUOTE: the longitudinal (P) wave propagates at the isotropic-solid speed $c_L = \sqrt{(K_{vac} + \tfrac{4}{3}G_{vac})/\rho_{bulk}} = \sqrt{\tfrac{10}{3}}\,c \approx 1.83c$
%          STAMPED AT: :11
%          RATIONALE: Objectives box asserts the propagating P-wave outright, plus sqrt(2)c as 'the
%          dilatational compression speed'; also still carries the stale 'canonical derivation in Vol 2 Ch 7'
%          for nu=2/7 that vol1 ch04 already corrected to GR-imported (comment twin at :10).
%          ALSO COVERED (named in the audited row, not separately stamped): :10
% --------------------------------------------------------------------------
% --------------------------------------------------------------------------
% R40 batch-2a --- NEEDS-RE-DERIVATION status note (2026-08-11)
% --------------------------------------------------------------------------
% CLASS: status demotion under R40. Mints no clm-/def-/exp-/sup-/ilk-, moves no solidity number,
% adjudicates no channel and opens no fork. Every byte of each demoted claim is preserved; the
% stamped line gains a status marker only.
%
% THE ARC, IN FOUR CLAUSES (R40's header form; clause 4 points at the LANDED artifact, not at a
% ruling record). (1) The kill fired --- the walk-back that closed the bulk radiative-port
% reading. (2) The premise localized to the imported K = 2G elastic modulus: the compressible
% far-field branch was minted by a GR-imported modulus, not forced by the axioms. (3) The axioms
% underdetermine the bulk sector --- the flat-direction finding: the written action conserves the
% Gauss function pointwise and never fixes its value. (4) THE REPLACEMENT IS THE LANDED RATIFIED
% BOUND-SECTOR LAW --- AXIOM 5, SUBSTRATE DC BIAS, clauses S (deposit), G (bias coupling / bridge)
% and Q (quiescence), canonical at manuscript/common_equations/eq_axiom_5.tex with its register
% entry in manuscript/ave-kb/common/axiom-register.md. Under clause G the A1 / bulk slot is a
% BOUND RESPONSE --- u_0 = -A_g grad(eps_11), mechanism gloss BACK-REACTION --- with no
% independent propagating branch, no port and zero longitudinal characteristic speed. A bulk wave
% speed, a bulk radiative port, a bulk band-branch and a bulk transit clock therefore have no
% referent. A_g (the bias-coupling area) is an UNVALUED-RATIFIED-CONSTANT per R48
% (manuscript/ave-kb/common/interlock-register.md): it is not valued here or anywhere, and THE
% CALIBRATION COUNT STAYS 3.
%
% STANDING NAMED-OPEN DEBT (the honesty rider). The ratified axiom does NOT discharge everything.
% THE BIAS PROPAGATION THEOREM is Axiom 5's standing named-open debt, stated by the axiom's own
% phase-structure paragraph, clause (c1): clause G's elliptic law is the static abstraction of
% underived finite-speed bias dynamics, and the (u,pi) no-signalling theorem does NOT cover the
% bias read --- the bias's finite propagation speed is owed, not held. Every row tagged BIAS-DEBT
% below re-derives against the ratified axiom WITH THAT DEBT STANDING, never against a closed
% replacement.
%
% VOCABULARY. Canonical nouns authored here: the bound response (u_0), the bias (eps_11), the DC
% operating point / quiescent point (Q-point); back-reaction is the mechanism gloss. 'dress',
% 'grade' as eps_11's canonical noun, and 'halo' for the physics (the physics noun is the
% near-field store / added-mass) are RETIRED by R50; 'retardation' is retired by R49(b) in favour
% of propagation delay / finite propagation speed. Corpus text quoted below is byte-exact and is
% never reworded.
%
% ROWS CARRIED IN THIS FILE (verbatim quote + verbatim banked rationale):
%
%   :163  family: K-backed stability (Cauchy implosion)  [BIAS-DEBT]
%        QUOTE (content verified at HEAD; markup-reduced from the banked audit): This positive bulk resistance guarantees that the spatial substrate
%        is incompressible and thermodynamically stable against gravitational collapse.
%       STAMPED AT: :163
%        RATIONALE (verbatim, research/drivers/r40_sweep_worklist_verified.json): K=2G as stability
%        guarantor — the prereg's named jeopardy family (cauchy-implosion-resolution twin); the carve
%        supplies stability kinematically (no compression DOF to implode); note the one-import-two-duties
%        symptom: 'incompressible' bought with the same K that mints a compressible P-branch elsewhere.
%        RESOLUTION: the demoted carrier is the propagating A1 / bulk branch; under Axiom 5 clause G that
%        slot is the bound response, so the re-derivation must be re-posed on the bound-sector
%        constitutive law (bias eps_11, bound response u_0, mechanism gloss back-reaction) rather than on
%        a compression wave. BIAS-DEBT: this row turns on finite-speed bias dynamics, so the resolution
%        is the ratified axiom WITH THE BIAS PROPAGATION THEOREM STANDING (clause (c1)) --- the
%        replacement is owed, not held.
%
% RECORDS: ruling R40 (the demotion sweep); the banked worklist
% research/drivers/r40_sweep_worklist_verified.json; the batch-0 scope verification and batch-1
% execution records in _orchestration/; this batch's record
% _orchestration/2026-08-12_r40-sweep-batch2a.md.
% --------------------------------------------------------------------------
